\section{Conclusão}
\label{sec:conclusao}

Este trabalho replicou, para a linha 4110 de Belo Horizonte, a análise de previsibilidade de
tempos de deslocamento com Modelos Mistos proposta por \cite{courb} para
Curitiba. Com 393 viagens observadas em 16 dias, o intercepto aleatório por trecho explicou
63\% da variância, e o modelo com erro AR(1) reduziu o erro de previsão fora da amostra de
62~s para 35~s por trecho. O padrão foi estável entre as duas semanas analisadas, o que indica
que a metodologia é transferível para outra cidade e outra fonte de dados.

A extensão proposta mostrou que as viagens duraram em média 70~minutos, contra 53~minutos no
quadro de horários, e que o atraso cresce 0,92~minuto por quilômetro percorrido. Os trechos
mais lentos, na região da BR-356 e no centro, concentram boa parte desse acúmulo. Esses
resultados podem subsidiar a revisão do quadro de horários da linha e a priorização de
intervenções nos trechos críticos. Por outro lado, os dados não permitiram estimar com
precisão o efeito do pico da tarde, e não foi observada piora entre as semanas.

\subsection{Limitações}

\begin{itemize}
  \item A coleta capturou apenas o sentido Belvedere $\rightarrow$ Dom Cabral; os resultados
        não podem ser estendidos ao sentido contrário.
  \item Foram 16 dias de coleta, com lacunas, e poucas viagens completas no início da manhã,
        no pico da tarde e aos domingos.
  \item Mais da metade das viagens com registro foi descartada por cobertura insuficiente,
        o que pode introduzir viés se as viagens descartadas tiverem comportamento diferente.
  \item Os horários programados nas paradas intermediárias são estimativas do GTFS, e não
        horários fiscalizados.
  \item O atraso é medido em relação ao \textit{trip\_id} informado pela API; se um veículo
        for rotulado com a viagem errada, o atraso dessa viagem fica distorcido. Os tempos de
        deslocamento entre paradas, usados na réplica, não dependem desse rótulo.
\end{itemize}

\subsection{Trabalhos futuros}

\begin{itemize}
  \item Coletar os dois sentidos da linha por um período mais longo, com monitoramento da
        coleta para reduzir as lacunas, sobretudo no pico.
  \item Incluir variáveis explicativas nos modelos, como chuva, distância e número de
        semáforos de cada trecho, e testar inclinações aleatórias por faixa horária.
  \item Aplicar a metodologia a outras linhas de Belo Horizonte, incluindo linhas do MOVE, e
        comparar a previsibilidade entre corredores.
  \item Usar os modelos para gerar previsões de chegada em tempo real e compará-las com as
        previsões disponibilizadas aos passageiros.
\end{itemize}
